% Generated draft: requires mathematical and source review.
\documentclass[11pt]{article}
\usepackage[margin=25mm]{geometry}
\usepackage{fontspec}
\setmainfont{Latin Modern Roman}
\usepackage{amsmath,amssymb,mathtools,unicode-math}
\setmathfont{Latin Modern Math}
\usepackage{enumitem,needspace,xurl,hyperref}
\setlength{\parindent}{0pt}
\setlength{\parskip}{5pt}
\setlength{\emergencystretch}{3em}
\newcommand{\N}{\mathbb{N}}
\newcommand{\Z}{\mathbb{Z}}
\newcommand{\Q}{\mathbb{Q}}
\newcommand{\R}{\mathbb{R}}
\newcommand{\C}{\mathbb{C}}
\newcommand{\F}{\mathbb{F}}
\newcommand{\A}{\mathbb{A}}
\newcommand{\PP}{\mathbb{P}}
\newcommand{\E}{\mathbb{E}}
\newcommand{\eps}{\varepsilon}
\newcommand{\dd}{\,\mathrm{d}}
\begin{document}
\section*{UnsolvedMath ID: 30002655}
\textbf{Source problem code:} catalogue entry
\subsection{Definitions and mathematical statement}
A subset of an ordered set satisfies the descending chain condition (DCC) if it contains no infinite strictly decreasing sequence. Fix a dimension \(n\), a DCC set \(I\subseteq[0,1]\), and a DCC set \(J\subseteq\R_{\geq0}\).

In the usual birational formulation, a pair \((X,\Delta)\) consists of a normal variety \(X\) and an effective real Weil divisor \(\Delta\) such that \(K_X+\Delta\) is real Cartier. The coefficient condition means that every nonzero coefficient of \(\Delta\) belongs to \(I\). For a prime divisor \(E\) on a normal birational model \(f:Y\to X\), its log discrepancy is
\[
a(E;X,\Delta)=1+\operatorname{coeff}_E\bigl(K_Y-f^*(K_X+\Delta)\bigr),
\]
with compatible canonical divisors. The pair is log canonical if \(a(E;X,\Delta)\geq0\) for every prime divisor over \(X\).

\textbf{Unresolved qualification.} The supplied record uses ``the log discrepancy'' of a pair as a single number without defining it. A divisorial log discrepancy requires a choice of \(E\); a minimal log discrepancy requires an infimum over specified divisors, usually with a specified center or locus. These are different formulations. The record also does not specify the base field or its characteristic. The definitions above explain standard terminology but do not resolve these omissions.

\textbf{Statement template from the supplied record.} Once the intended class of varieties and the intended scalar invariant \(A(X,\Delta)\) denoted by ``log discrepancy'' have been specified, the requested assertion is
\[
\begin{gathered}
\exists\,J_0\subseteq J,\quad |J_0|<\infty,\quad
\forall\,(X,\Delta),\\
\bigl[\dim X=n,\ (X,\Delta)\text{ log canonical},\
\operatorname{coeff}(\Delta)\subseteq I,\ A(X,\Delta)\in J\bigr]
\Longrightarrow A(X,\Delta)\in J_0.
\end{gathered}
\]
Here \(J_0\) is uniform over all pairs under consideration, for the fixed data \(n,I,J\). If the intended invariant additionally depends on a divisor, point, or locus, those data and their quantifiers must be supplied. Thus this is a conditional transcription of the recorded question, not a fully specified conjecture.
\subsection{Short English statement}
The recorded question asks whether, in fixed dimension and with boundary coefficients in a fixed DCC set, only finitely many permitted log discrepancy values can lie in any fixed DCC set \(J\). The meaning of the numerical invariant, its required center or divisor data, and the underlying field remain unspecified. The supplied record attributes the conjecture to Shokurov and labels it open; this draft does not independently verify that attribution or status.
\subsection{Sources}
\begin{itemize}
\item Local UnsolvedMath registry: \texttt{data/unsolvedmath/problems.json}, supplied record with permanent ID \(30002655\), exact source problem code \texttt{catalogue entry}, title ``Ascending Chain Condition for Log Discrepancies.'' The record belongs to the Oberwolfach Reports Open Problems collection and gives proposed year \(2014\); no individual proposer is recorded.
\item Supplied report citation: \emph{Komplexe Analysis (2015)}, Oberwolfach Reports, pp.\ 2177--2234, DOI \texttt{10.4171/owr/2014/39}. Original Oberwolfach report: \url{https://doi.org/10.4171/owr/2014/39}. The background attributes the conjecture to Shokurov. The relationship between the citation's year \(2015\) and the record's \(2014\) dating is unresolved here.
\item Supplied literature reference, labeled ``Kawakita, smooth-threefold ACC result'': \url{https://doi.org/10.1112/jlms.70510}. No theorem statement or detailed hypotheses are supplied.
\item Literature assessment dated exactly \texttt{2026-08-21}: the supplied curation labels the problem open, states that the extracted item names Shokurov's conjecture without resolving it, and reports results only in special dimensions or singularity classes. This is supplied dated curation, not an independently verified literature assessment. Its own qualification requests independent verification of time-sensitive status and cited claims.
\item The registry labels the statement \texttt{corrected\_verified} and reports that it was checked against the official report to restore hypotheses and notation. That verification claim is supplied metadata; the report and literature references have not been independently consulted for this draft.
\end{itemize}
\end{document}
